\documentclass{summarysheet}

\begin{document}

\maketitle{10}{Rotation}


\begin{multicols}{2}

\begin{topicbox}{The Rigid Body Model}

\noindent A \emph{rigid body} is an extended object whose size and shape don't change.  Rigid bodies can't be stretched, compressed, or deformed; all points on the rigid body have the same angular velocity and angular acceleration; and rigid bodies can have translational motion, rotational motion, or a combination of both.

\end{topicbox}


\begin{topicbox}{Moment of Inertia}

\noindent The \emph{moment of inertia} is the rotational equivalent of mass, and can be calculated for $N$ particles of mass $m_i$ by
\begin{eqbox}
I = \sum_i^N  m_i r_i^2,
\end{eqbox}
\noindent where $r_i$ is the distance from the axis of rotation.

%For objects that are continuous, the moment of inertia is calculated from
%\[
%I = \int r^2 \, dm,
%\]
%although it's easier to just look up the value based on its shape and rotation axis:

For some objects, you can just look up the moment of inertia:
\begin{center}
\begin{tabular}{ll|ll}
\textbf{Shape, axis} & \textbf{$I$} &  \textbf{Shape, axis} & \textbf{$I$}\\
\hline \\
Thin rod, centre & $(1/12) ML^2$ & Thin rod, end & $(1/3) ML^2$ \\
Plane, centre & $(1/12) Ma^2$ & Plane, edge & $(1/3) Ma^2$ \\
Disk & $(1/2) MR^2$ & Hoop & $MR^2$ \\
Solid sphere & $(2/5) MR^2$ & Spherical shell &$(2/3) MR^2$ 
\end{tabular}
\end{center}

The parallel axis theorem allows you to find the moment of inertia about an axis parallel to the centre of mass axis:
\[
I = I_\text{cm} + Md^2,
\]
where $d$ is the distance from the centre of mass.

\end{topicbox}



\begin{topicbox}{The Cross Product}

\noindent Two vectors can be multiplied together to get another vector through the \emph{cross-product}:
\begin{eqbox}
\vec{A} \times \vec{B} = (AB \sin\alpha, \text{direction by the right hand rule}).
\end{eqbox}

%\begin{center}
%\includegraphics[scale=0.7]{fig_cross.pdf}
%\end{center}

\end{topicbox}



\begin{topicbox}{Torque}

\noindent Torque is the rotational equivalent of force, and is defined to be

%\begin{multicols}{2}
\begin{eqbox}
\vec{\tau} = \vec{r} \times \vec{F},
\end{eqbox}
\noindent where $\vec{r}$ is the position of where the force is applied measured from the rotation axis.

Note that \emph{gravitational torque} can be found by treating the object as if all its mass were concentrated at the centre of mass:
\[
\tau_\text{grav} = -M g x_\text{cm}.
\]

%\begin{center}
%\includegraphics[scale=0.7]{fig_torque.pdf}
%\end{center}

%\end{multicols}

\end{topicbox}



\begin{topicbox}{Rotational Kinetic Energy}

\noindent A rigid body object rotating with angular speed $\omega$ has a rotational kinetic energy given by
\begin{eqbox}
K_\text{rot} = \frac{1}{2} I \omega^2.
\end{eqbox}

\end{topicbox}



\begin{topicbox}{Centre of Mass}

\noindent  An unconstrained object of mass $M$ with no net force will rotate about its \emph{centre of mass}: the only point that remains motionless.  The centre of mass can be found by
\[
x_\text{cm} = \frac{1}{M} \int x\, dm \quad \text{and} \quad y_\text{cm} = \frac{1}{M} \int y\, dm,
\]
or, if the object consists of $N$ particles of mass $m_i$, 
\begin{eqbox}
x_\text{cm} = \frac{1}{M} \sum_i^N m_i x_i \quad \text{and} \quad y_\text{cm} = \frac{1}{M} \sum_i^N m_i y_i.
\end{eqbox}

\end{topicbox}

\begin{topicbox}{Newton's Second Law for Rotation}

\noindent A net torque causes an angular acceleration given by Newton's second law,
\begin{eqbox}
\alpha = \frac{\tau_\text{net}}{I}.
\end{eqbox}

If an object is stationary, both the \emph{net force} and \emph{net torque} are zero.  This is called \emph{static equilibrium}:
\[
\text{Static equilibrium:} \quad \vec{F}_\text{net} = \vec{0} \quad \text{and} \quad \vec{\tau}_\text{net} = \vec{0}.
\]

\end{topicbox}













\begin{topicbox}{Rolling}

\noindent Rolling consists of both \emph{translational} and \emph{rotational} motion; these motions are linked by the \emph{rolling constraint},
\[
v_\text{cm} = R\omega.
\]
The rolling of a rigid body can ve described as a translation of the centre of mass plus a rotation about the centre of mass.  The total kinetic energy of rolling is
\[
K = K_\text{rot} + K_\text{cm} = \frac{1}{2} I_\text{cm} \omega^2 + \frac{1}{2}M v_\text{cm}^2.
\]
\end{topicbox}


\begin{topicbox}{Angular Momentum}

\noindent The \emph{angular momentum} of a \emph{particle} is defined as 
\begin{eqbox}
\vec{L} = \vec{r} \times \vec{p}.
\end{eqbox}
\noindent In the case that a \emph{rigid body} is rotating about a fixed axis or an axis of symmetry, its angular momentum is 
\[
\vec{L} = I \vec{\omega}.
\]

We can write Newton's second law for rotation in terms of angular momentum:
\[
\frac{d\vec{L}}{dt} = \vec{\tau}_\text{net},
\]
so that a net torque changes the angular momentum of an object.

If the system is \emph{isolated}, so that $\tau_\text{net} = 0$, angular momentum is conserved:
\begin{eqbox}
\vec{L}_i = \vec{L}_f.
\end{eqbox}

\end{topicbox}


\end{multicols}





\makebanner{Mechanics}

\end{document}




